
\chapter{Практическая часть}

\section{Программная реализация}

В листингах \ref{lst:main.m} и \ref{lst:hungarian.m} приведен текст программы.

\includelistingpretty 
    {main.m} % Имя файла с расширением (файл должен быть расположен в директории inc/lst/) 
    {matlab} % Язык программирования (необязательный аргумент) 
    {main.m} % Подпись листинга

\includelistingpretty 
    {hungarian.m} % Имя файла с расширением (файл должен быть расположен в директории inc/lst/) 
    {matlab} % Язык программирования (необязательный аргумент) 
    {hungarian.m} % Подпись листинга

\section{Результаты решений}

\subsection{Задача минимизации}

\includelistingpretty 
    {out_min.m} % Имя файла с расширением (файл должен быть расположен в директории inc/lst/) 
    {matlab} % Язык программирования (необязательный аргумент) 
    {Вывод программы при задаче минимизации} % Подпись листинга


% Исходная матрица стоимостей:

% \begin{equation}
%     C = 
% \begin{bmatrix}
% 3 & 5 & 2 & 4 & 8\\
% 10 & 10 & 4 & 3 & 6\\
% 5 & 6 & 9 & 8 & 3\\
% 6 & 2 & 5 & 8 & 4\\
% 5 & 4 & 8 & 9 & 3
% \end{bmatrix}
% \end{equation}

% После вычитания минимальных элементов по столбцам получаем матрицу представленную в \ref{equ:matrix1}:

% \begin{equation}\label{equ:matrix1}
%     C \rightarrow
% \begin{bmatrix}
% 0 & 3 & 0 & 1 & 5\\
% 7 & 8 & 2 & 0 & 3\\
% 2 & 4 & 7 & 5 & 0\\
% 3 & 0 & 3 & 5 & 1\\
% 2 & 2 & 6 & 6 & 0
% \end{bmatrix}
% \end{equation}


% После вычитания минимальных элементов по строкам получаем матрицу предстваленную в \ref{equ:matrix2}:

% \begin{equation}\label{equ:matrix2}
%     C \rightarrow
% \begin{bmatrix}
% 0 & 3 & 0 & 1 & 5\\
% 7 & 8 & 2 & 0 & 3\\
% 2 & 4 & 7 & 5 & 0\\
% 3 & 0 & 3 & 5 & 1\\
% 2 & 2 & 6 & 6 & 0
% \end{bmatrix}
% \end{equation}

% Строим первичную СНН: первый найденный в столбце 0, в одной строке с которым нет элемента $0^*$, отмечаем как $0^*$. Получаем матрицу представленную в \ref{equ:matrix3}.

% \begin{equation}\label{equ:matrix3}
%     C \rightarrow
% \begin{bmatrix}
% 0^* & 3 & 0 & 1 & 5\\
% 7 & 8 & 2 & 0^* & 3\\
% 2 & 4 & 7 & 5 & 0^*\\
% 3 & 0^* & 3 & 5 & 1\\
% 2 & 2 & 6 & 6 & 0
% \end{bmatrix}
% \end{equation}


% Число независимых нулей $k = 4 < n = 5$, следовательно необходимо улучшение СНН.

% Отмечаем столбцы, где есть $0^*$ знаком $+$. Находим 0 среди невыделенных элементов и отмечаем его $0'$. Результат представлен в \ref{equ:matrix4}:

% \begin{equation}\label{equ:matrix4}
%     C \rightarrow
% \begin{bmatrix}
% 0^* & 3 & 0' & 1 & 5\\
% 7 & 8 & 2 & 0^* & 3\\
% 2 & 4 & 7 & 5 & 0^*\\
% 3 & 0^* & 3 & 5 & 1\\
% 2 & 2 & 6 & 6 & 0
% \end{bmatrix}
% \end{equation}

% \begin{equation*}
% \begin{matrix}
%     & "" & "" & + & + & - & + & +
% \end{matrix}
% \end{equation*}

% Далее если в одной строке с $0'$ стоит $0^*$, снимаем выделение со столбца с найденным $0^*$ и отмечаем строку с текущим $0'$, как представлено в выражении \ref{equ:matrix5}.

% \begin{equation}\label{equ:matrix5}
%     C \rightarrow
% \begin{bmatrix}
% 0^* & 3 & 0' & 1 & 5\\
% 7 & 8 & 2 & 0^* & 3\\
% 2 & 4 & 7 & 5 & 0^*\\
% 3 & 0^* & 3 & 5 & 1\\
% 2 & 2 & 6 & 6 & 0
% \end{bmatrix}
% \begin{matrix}
% + \\
% - \\
% - \\
% - \\
% -
% \end{matrix}
% \end{equation}

% \begin{equation*}
% \begin{matrix}
%     & "" & "" & - & + & - & + & +
% \end{matrix}
% \end{equation*}

% Среди невыделенных элементов нет нулей. Выбираем наименьший элемент $h = 2$ среди невыделенных. Вычитаем h из элементов в невыделенных столбцах, прибавляем h к элементам в выделенных строках.После перерасчёта матрицы $h$ получаем результат представленный в \ref{equ:matrix6}:

% \begin{equation}\label{equ:matrix6}
%     C \rightarrow
% \begin{bmatrix}
% 0^* & 5 & 0' & 3 & 7\\
% 5 & 8 & 0 & 0^* & 3\\
% 0 & 4 & 5 & 5 & 0^*\\
% 1 & 0^* & 1 & 5 & 1\\
% 0 & 2 & 4 & 6 & 0
% \end{bmatrix}
% \begin{matrix}
% + \\
% - \\
% - \\
% - \\
% -
% \end{matrix}
% \end{equation}

% \begin{equation*}
% \begin{matrix}
%     & "" & "" & - & + & - & + & +
% \end{matrix}
% \end{equation*}


% Находим 0 среди невыделенных элементов и отмечаем его $0'$. Результат представлен в \ref{equ:matrix7}:

% \begin{equation}\label{equ:matrix7}
%     C \rightarrow
% \begin{bmatrix}
% 0^* & 5 & 0' & 3 & 7\\
% 5 & 8 & 0 & 0^* & 3\\
% 0 & 4 & 5 & 5 & 0^*\\
% 1 & 0^* & 1 & 5 & 1\\
% 0' & 2 & 4 & 6 & 0
% \end{bmatrix}
% \begin{matrix}
% + \\
% - \\
% - \\
% - \\
% -
% \end{matrix}
% \end{equation}

% \begin{equation*}
% \begin{matrix}
%     & "" & "" & - & + & - & + & +
% \end{matrix}
% \end{equation*}

% C найденным $0'$ в одной строке нет $0^*$. Можно строить L-цепочку. В получившейся L-цепочке заменяем $0^*$ на $0$, а $0'$ на $0^*$.Результат построения представлен в \ref{equ:matrix8}.

% \begin{equation}\label{equ:matrix8}
%     C \rightarrow
% \begin{bmatrix}
% 0 & 5 & 0^* & 3 & 7\\
% 5 & 8 & 0 & 0^* & 3\\
% 0 & 4 & 5 & 5 & 0^*\\
% 1 & 0^* & 1 & 5 & 1\\
% 0^* & 2 & 4 & 6 & 0
% \end{bmatrix}
% \end{equation}

% Число независимых нулей в матрице представленной в \ref{equ:matrix8} $k = 5 = n$, следовательно можно выписать оптимальную матрицу назначений, представленную в \ref{equ:matrix9}.

% \begin{equation}\label{equ:matrix9}
%     X_\text{opt} =
% \begin{bmatrix}
% 0 & 0 & 1 & 0 & 0\\
% 0 & 0 & 0 & 1 & 0\\
% 0 & 0 & 0 & 0 & 1\\
% 0 & 1 & 0 & 0 & 0\\
% 1 & 0 & 0 & 0 & 0
% \end{bmatrix}
% \end{equation}

% Таким образом минимальная суммарная стоимость может быть найдена из выражения, представленного в \ref{equ:eq1}:

% \begin{equation}\label{equ:eq1}
%     f_\text{opt} = 2 + 3 + 3 + 5 + 2 = 15
% \end{equation}

\subsection{Задача максимизации}

\includelistingpretty 
    {out_max.m} % Имя файла с расширением (файл должен быть расположен в директории inc/lst/) 
    {matlab} % Язык программирования (необязательный аргумент) 
    {Вывод программы при задаче максимизации} % Подпись листинга

% Исходная матрица стоимостей:

% \begin{equation}
%     C = 
% \begin{bmatrix}
% 3 & 5 & 2 & 4 & 8\\
% 10 & 10 & 4 & 3 & 6\\
% 5 & 6 & 9 & 8 & 3\\
% 6 & 2 & 5 & 8 & 4\\
% 5 & 4 & 8 & 9 & 3
% \end{bmatrix}
% \end{equation}

% Для решения задачи максимизации необходимо преобразовать матрицу: найти максимальный элемент исходной матрицы, вычесть его из всех элементов. Получившуюся матрицу умножить на (-1). Результат преобразования представлан на \ref{equ:matrix10}

% \begin{equation}\label{equ:matrix10}
%     C \rightarrow 
% \begin{bmatrix}
% 7 & 5 & 8 & 6 & 2 \\
% 0 & 0 & 6 & 7 & 4 \\
% 5 & 4 & 1 & 2 & 7 \\
% 4 & 8 & 5 & 2 & 6 \\
% 5 & 6 & 2 & 1 & 7
% \end{bmatrix}
% \end{equation}

% Дальнейшее решение аналогично решению задачи минимизации. После вычитания минимальных элементов по столбцам получаем матрицу представленную в \ref{equ:matrix11}:

% \begin{equation}\label{equ:matrix11}
%     C \rightarrow 
% \begin{bmatrix}
% 7 & 5 & 7 & 5 & 0 \\
% 0 & 0 & 5 & 6 & 2 \\
% 5 & 4 & 0 & 1 & 5 \\
% 4 & 8 & 4 & 1 & 4 \\
% 5 & 6 & 1 & 0 & 5
% \end{bmatrix}
% \end{equation}

% После вычитания минимальных элементов по строкам получаем матрицу представленную в \ref{equ:matrix12}:

% \begin{equation}\label{equ:matrix12}
%     C \rightarrow 
% \begin{bmatrix}
% 7 & 5 & 7 & 5 & 0 \\
% 0 & 0 & 5 & 6 & 2 \\
% 5 & 4 & 0 & 1 & 5 \\
% 3 & 7 & 3 & 0 & 3 \\
% 5 & 6 & 1 & 0 & 5
% \end{bmatrix}
% \end{equation}

% Строим первичную СНН: первый найденный в столбце 0, в одной строке с которым нет элемента $0^*$, отмечаем как $0^*$. Получаем матрицу представленную в \ref{equ:matrix13}.

% \begin{equation}\label{equ:matrix13}
%     C \rightarrow 
% \begin{bmatrix}
% 7 & 5 & 7 & 5 & 0^* \\
% 0^* & 0 & 5 & 6 & 2 \\
% 5 & 4 & 0^* & 1 & 5 \\
% 3 & 7 & 3 & 0^* & 3 \\
% 5 & 6 & 1 & 0 & 5
% \end{bmatrix}
% \end{equation}

% Число независимых нулей $k = 4 < n = 5$, следовательно необходимо улучшение СНН.

% Отмечаем столбцы, где есть $0^*$ знаком $+$. Находим 0 среди невыделенных элементов и отмечаем его $0'$. Результат представлен в \ref{equ:matrix14}:

% \begin{equation}\label{equ:matrix14}
%     C \rightarrow 
% \begin{bmatrix}
% 7 & 5 & 7 & 5 & 0^* \\
% 0^* & 0' & 5 & 6 & 2 \\
% 5 & 4 & 0^* & 1 & 5 \\
% 3 & 7 & 3 & 0^* & 3 \\
% 5 & 6 & 1 & 0 & 5
% \end{bmatrix}
% \end{equation}

% \begin{equation*}
% \begin{matrix}
%     & "" & "" & + & - & + & + & +
% \end{matrix}
% \end{equation*}

% Далее если в одной строке с $0'$ стоит $0^*$, снимаем выделение со столбца с найденным $0^*$ и отмечаем строку с текущим $0'$, как представлено в выражении \ref{equ:matrix15}.

% \begin{equation}\label{equ:matrix15}
%     C \rightarrow
% \begin{bmatrix}
% 7 & 5 & 7 & 5 & 0^* \\
% 0^* & 0' & 5 & 6 & 2 \\
% 5 & 4 & 0^* & 1 & 5 \\
% 3 & 7 & 3 & 0^* & 3 \\
% 5 & 6 & 1 & 0 & 5
% \end{bmatrix}
% \begin{matrix}
% - \\
% + \\
% - \\
% - \\
% -
% \end{matrix}
% \end{equation}

% \begin{equation*}
% \begin{matrix}
%     & "" & "" & - & - & + & + & +
% \end{matrix}
% \end{equation*}

% Среди невыделенных элементов нет нулей. Выбираем наименьший элемент $h = 3$ среди невыделенных. Вычитаем h из элементов в невыделенных столбцах, прибавляем h к элементам в выделенных строках.После перерасчёта матрицы $h$ получаем результат представленный в \ref{equ:matrix16}:

% \begin{equation}\label{equ:matrix16}
%     C \rightarrow
% \begin{bmatrix}
% 4 & 2 & 7 & 5 & 0^* \\
% 0^* & 0' & 8 & 9 & 5 \\
% 2 & 1 & 0^* & 1 & 5 \\
% 0 & 4 & 3 & 0^* & 3 \\
% 2 & 3 & 1 & 0 & 5
% \end{bmatrix}
% \begin{matrix}
% - \\
% + \\
% - \\
% - \\
% -
% \end{matrix}
% \end{equation}

% \begin{equation*}
% \begin{matrix}
%     & "" & "" & - & - & + & + & +
% \end{matrix}
% \end{equation*}

% Находим 0 среди невыделенных элементов и отмечаем его $0'$. Результат представлен в \ref{equ:matrix17}:

% \begin{equation}\label{equ:matrix17}
%     C \rightarrow
% \begin{bmatrix}
% 4 & 2 & 7 & 5 & 0^* \\
% 0^* & 0' & 8 & 9 & 5 \\
% 2 & 1 & 0^* & 1 & 5 \\
% 0' & 4 & 3 & 0^* & 3 \\
% 2 & 3 & 1 & 0 & 5
% \end{bmatrix}
% \begin{matrix}
% - \\
% + \\
% - \\
% - \\
% -
% \end{matrix}
% \end{equation}

% \begin{equation*}
% \begin{matrix}
%     & "" & "" & - & - & + & + & +
% \end{matrix}
% \end{equation*}

% Далее если в одной строке с $0'$ стоит $0^*$, снимаем выделение со столбца с найденным $0^*$ и отмечаем строку с текущим $0'$, как представлено в выражении \ref{equ:matrix18}.

% \begin{equation}\label{equ:matrix18}
%     C \rightarrow
% \begin{bmatrix}
% 4 & 2 & 7 & 5 & 0^* \\
% 0^* & 0' & 8 & 9 & 5 \\
% 2 & 1 & 0^* & 1 & 5 \\
% 0' & 4 & 3 & 0^* & 3 \\
% 2 & 3 & 1 & 0 & 5
% \end{bmatrix}
% \begin{matrix}
% - \\
% + \\
% - \\
% + \\
% -
% \end{matrix}
% \end{equation}

% \begin{equation*}
% \begin{matrix}
%     & "" & "" & - & - & + & - & +
% \end{matrix}
% \end{equation*}

% Находим 0 среди невыделенных элементов и отмечаем его $0'$. Результат представлен в \ref{equ:matrix19}:

% \begin{equation}\label{equ:matrix19}
%     C \rightarrow
% \begin{bmatrix}
% 4 & 2 & 7 & 5 & 0^* \\
% 0^* & 0' & 8 & 9 & 5 \\
% 2 & 1 & 0^* & 1 & 5 \\
% 0' & 4 & 3 & 0^* & 3 \\
% 2 & 3 & 1 & 0' & 5
% \end{bmatrix}
% \begin{matrix}
% - \\
% + \\
% - \\
% + \\
% -
% \end{matrix}
% \end{equation}

% \begin{equation*}
% \begin{matrix}
%     & "" & "" & - & - & + & - & +
% \end{matrix}
% \end{equation*}

% C найденным $0'$ в одной строке нет $0^*$. Можно строить L-цепочку. В получившейся L-цепочке заменяем $0^*$ на $0$, а $0'$ на $0^*$.Результат построения представлен в \ref{equ:matrix20}.

% \begin{equation}\label{equ:matrix20}
%     C \rightarrow
% \begin{bmatrix}
% 4 & 2 & 7 & 5 & 0^* \\
% 0 & 0^* & 8 & 9 & 5 \\
% 2 & 1 & 0^* & 1 & 5 \\
% 0^* & 4 & 3 & 0 & 3 \\
% 2 & 3 & 1 & 0^* & 5
% \end{bmatrix}
% \end{equation}

% Число независимых нулей в матрице представленной в \ref{equ:matrix20} $k = 5 = n$, следовательно можно выписать оптимальную матрицу назначений, представленную в \ref{equ:matrix21}.

% \begin{equation}\label{equ:matrix21}
%     X_\text{opt} =
% \begin{bmatrix}
% 0 & 0 & 0 & 0 & 1 \\
% 0 & 1 & 0 & 0 & 0 \\
% 0 & 0 & 1 & 0 & 0 \\
% 1 & 0 & 0 & 0 & 0 \\
% 0 & 0 & 0 & 1 & 0
% \end{bmatrix}
% \end{equation}

% Таким образом максимальная суммарная стоимость может быть найдена из выражения, представленного в \ref{equ:eq2}:

% \begin{equation}\label{equ:eq2}
%     f_\text{opt} = 6 + 10 + 9 + 9 + 8 = 42 
% \end{equation}

\chapter*{ЗАКЛЮЧЕНИЕ}

В результате выполнения лабораторной работы был реализован венгерский метод решения задачи о назначениях для задач минимизации и максимизации. В ходе выполнения лабораторной работы была написана
программа на языке MatLab, позволяющая применять венгерский метод решения задачи о назначениях.